\section{Conclusion}
We presented SPRM, a Shapley-inspired process reward modeling framework that
constructs step-level supervision directly from outcome-only reasoning
trajectories.
By combining local attribution, causal consistency checking, and a global
semantic prior, SPRM improves over comparable baseline PRMs on both
ProcessBench and PRMBench and also transfers well to Best-of-\(N\) reranking.
These results suggest that dense and reusable process supervision can be
derived from sparse final outcomes without relying on external step annotation
or judge models.

\paragraph{Limitations and Broader Impacts}
Our study is limited to outcome-labeled reasoning trajectories and benchmarks
centered on mathematical or structured reasoning, so performance in broader
domains remains to be validated.
The method also introduces extra computational overhead from attribution and
cross-trajectory retrieval, and its global prior may degrade when semantic
neighbors are sparse or distributionally mismatched.
On the positive side, reducing dependence on external judges may lower the cost
and bias of process supervision.
On the negative side, stronger reasoning reward models could be used to scale
automated reasoning systems in settings where incorrect but persuasive outputs
are harmful, so controlled release and careful evaluation remain important.